\documentclass[11pt,a4paper]{article}
\usepackage[utf8]{inputenc}
\usepackage[margin=0.8in]{geometry}
\usepackage{listings}
\usepackage{xcolor}
\usepackage{titlesec}
\usepackage{hyperref}
\usepackage{fancyhdr}

\definecolor{codebg}{RGB}{245,245,245}
\definecolor{codegreen}{RGB}{0,128,0}
\definecolor{codegray}{RGB}{128,128,128}
\definecolor{codeblue}{RGB}{0,0,180}
\definecolor{codepurple}{RGB}{128,0,128}

\lstset{
    language=C++,
    backgroundcolor=\color{codebg},
    basicstyle=\ttfamily\footnotesize,
    keywordstyle=\color{codeblue}\bfseries,
    commentstyle=\color{codegreen}\itshape,
    stringstyle=\color{codepurple},
    numberstyle=\tiny\color{codegray},
    numbers=left,
    numbersep=5pt,
    frame=single,
    breaklines=true,
    breakatwhitespace=false,
    tabsize=4,
    showstringspaces=false,
    captionpos=b,
    morekeywords={override,nullptr,string,shared_ptr,unique_ptr,make_shared,make_unique,
                  lock_guard,unique_lock,mutex,thread,vector,unordered_map,map,
                  chrono,steady_clock,time_point,atomic},
}

\pagestyle{fancy}
\fancyhf{}
\fancyhead[L]{\textbf{LLD --- Parking Lot}}
\fancyhead[R]{\thepage}

\title{\Huge\textbf{Parking Lot}\\[0.3em]\Large Low-Level Design --- C++ Concurrency}
\author{LLD Interview Preparation}
\date{}

\begin{document}
\maketitle
\tableofcontents
\newpage

% ============================================================
\section{File Structure}
% ============================================================
\begin{verbatim}
parking-lot/
  models.hpp       -- Vehicle, ParkingSpot, ParkingFloor, Ticket
  strategies.hpp   -- PaymentStrategy, RateStrategy
  managers.hpp     -- ParkingManager (Singleton)
  main.cpp         -- Demo: concurrent parking/unparking
\end{verbatim}

\textbf{Compilation:}
\begin{verbatim}
g++ -std=c++17 -pthread -o parking_lot.exe main.cpp
\end{verbatim}

\subsection{Concurrency Summary}
\begin{itemize}
    \item \textbf{Level 1:} \texttt{operationMtx} (ParkingManager) --- floors, activeTickets, ticketCounter
    \item \textbf{Level 2:} None (single manager lock suffices)
    \item \textbf{lock\_guard:} \texttt{parkVehicle}, \texttt{setRateStrategy}, \texttt{addFloor}
    \item \textbf{unique\_lock:} \texttt{unparkVehicle} (release before payment I/O)
\end{itemize}

\subsection{Smart Pointers}
\begin{itemize}
    \item \texttt{shared\_ptr<Ticket>} --- manager map + caller both hold ref
    \item \texttt{raw ptr (RateStrategy*)} --- non-owning, caller owns strategy
    \item \texttt{raw ptr (Vehicle*)} --- non-owning, caller owns vehicles
\end{itemize}

% ============================================================
\newpage
\section{models.hpp}
% ============================================================
\begin{lstlisting}[title={\texttt{models.hpp}}]
#ifndef MODELS_HPP
#define MODELS_HPP

#include<bits/stdc++.h>
#include<string>
#include<chrono>
using namespace std;

enum class VehicleType{
    BIKE,
    CAR,
    TRUCK
};

enum class SpotType{
    SMALL,
    MEDIUM,
    LARGE
};

enum class TicketStatus{
    ACTIVE,
    PAID
};

class Vehicle{
  private:
    string licensePlate;
    VehicleType type;
  public:
    Vehicle() : licensePlate(""), type(VehicleType::CAR) {}
    Vehicle(string plate, VehicleType vtype): licensePlate(plate), type(vtype) {}

    VehicleType getType() const { return type; }
    string getLicensePlate() const { return licensePlate; }
};

class ParkingSpot{
    private:
        string spotId;
        bool isOccupied;
        SpotType type;
        Vehicle* currentVehicle;
     
    public:
        ParkingSpot() : spotId(""), isOccupied(false), type(SpotType::SMALL), currentVehicle(nullptr) {}
        ParkingSpot(string id, SpotType stype)
            : spotId(id), isOccupied(false), type(stype), currentVehicle(nullptr) {}

        bool canFit(VehicleType vtype){
            if(isOccupied) return false;
            if((vtype == VehicleType::BIKE && type == SpotType::SMALL) ||
               (vtype == VehicleType::CAR && type == SpotType::MEDIUM) ||
               (vtype == VehicleType::TRUCK && type == SpotType::LARGE)){
                return true;
            }
            return false;
        }

        void park(Vehicle* v){
            this->currentVehicle = v;
            isOccupied = true;
        }

        void unpark(){
            this->currentVehicle = nullptr;
            isOccupied = false;
        }

        string getSpotId() const { return spotId; }
        bool getIsOccupied() const { return isOccupied; }
        SpotType getType() const { return type; }
}; 


class ParkingFloor{
    private:
    int floorNumber;
    vector<ParkingSpot> spots;

    public:
    ParkingFloor() : floorNumber(0) {}
    ParkingFloor(int num, vector<ParkingSpot> s) : floorNumber(num), spots(s) {}

    ParkingSpot* findAvailableSpot(VehicleType vtype){
        for(auto& spot : spots){
            if(spot.canFit(vtype)){
                return &spot;
            }
        }
        return nullptr;
    }

    int getFloorNumber() const { return floorNumber; }
};

class Ticket{
    private:
    string ticketId;
    Vehicle vehicle;
    ParkingSpot* spot;
    chrono::steady_clock::time_point entryTime;
    chrono::steady_clock::time_point exitTime;
    TicketStatus status;
    double amount;

    public:
    Ticket() : spot(nullptr), status(TicketStatus::ACTIVE), amount(0.0) {}
    Ticket(string id, Vehicle v, ParkingSpot* s, chrono::steady_clock::time_point entry)
        : ticketId(id), vehicle(v), spot(s), entryTime(entry),
          exitTime(), status(TicketStatus::ACTIVE), amount(0.0) {}

    string getTicketId() const { return ticketId; }
    Vehicle getVehicle() const { return vehicle; }
    ParkingSpot* getSpot() const { return spot; }
    TicketStatus getStatus() const { return status; }
    double getAmount() const { return amount; }
    chrono::steady_clock::time_point getEntryTime() const { return entryTime; }
    chrono::steady_clock::time_point getExitTime() const { return exitTime; }

    long long getDurationMinutes() const {
        auto duration = chrono::duration_cast<chrono::minutes>(exitTime - entryTime);
        return max(1LL, duration.count());
    }

    void setExitTime(chrono::steady_clock::time_point t) { exitTime = t; }
    void setStatus(TicketStatus s) { status = s; }
    void setAmount(double a) { amount = a; }
};

#endif
\end{lstlisting}

% ============================================================
\newpage
\section{strategies.hpp}
% ============================================================
\begin{lstlisting}[title={\texttt{strategies.hpp}}]
#ifndef STRATEGIES_HPP
#define STRATEGIES_HPP

#include<iostream>
using namespace std;


class PaymentStrategy{
    public:
    virtual bool pay(double amount) = 0;
    virtual ~PaymentStrategy() {}
};

class CashPayment: public PaymentStrategy {
   public:
   bool pay(double amount) override {
       cout << "Paid $" << amount << " in cash." << endl;
       return true;
   }
};

class CardPayment: public PaymentStrategy {
   public:
   bool pay(double amount) override {
       cout << "Paid $" << amount << " by card." << endl;
       return true;
   }
};

class RateStrategy{
    public:
    virtual double calculateRate(long long minutes) = 0;
    virtual ~RateStrategy() {}
};

class HourlyRateStrategy: public RateStrategy {
    double ratePerHour;
    public:
    HourlyRateStrategy(double rate = 10.0) : ratePerHour(rate) {}
    double calculateRate(long long minutes) override {
        double hours = (minutes + 59) / 60.0;
        double amount = hours * ratePerHour;
        cout << "Hourly rate: " << ratePerHour << "/hr x " << hours << " hrs = $" << amount << endl;
        return amount;
    }
};

class FlatRateStrategy: public RateStrategy {
    double flatRate;
    public:
    FlatRateStrategy(double rate = 50.0) : flatRate(rate) {}
    double calculateRate(long long minutes) override {
        cout << "Flat rate applied: $" << flatRate << " (duration: " << minutes << " min)" << endl;
        return flatRate;
    }
};

#endif
\end{lstlisting}

% ============================================================
\newpage
\section{managers.hpp}
% ============================================================
\begin{lstlisting}[title={\texttt{managers.hpp}}]
#ifndef MANAGERS_HPP
#define MANAGERS_HPP

#include<bits/stdc++.h>
#include<mutex>
#include<memory>
#include<algorithm>
#include<thread>
#include "models.hpp"
#include "strategies.hpp"
using namespace std;

// ============================================================
// SMART POINTER CHEAT SHEET:
// ----------------------------------------------------------
// unique_ptr  -> sole ownership, non-copyable, lightweight
//   Use when: one owner (e.g., singleton instance, owned children)
// shared_ptr  -> shared ownership, ref-counted, heavier
//   Use when: multiple owners need to keep object alive
//   (e.g., Ticket returned to caller but also stored in map)
// raw ptr     -> non-owning reference (observer)
//   Use when: you don't own the object (e.g., rateStrategy set from outside)
// ============================================================

// ============================================================
// LOCK CHEAT SHEET:
// ----------------------------------------------------------
// lock_guard<mutex>   -> simple RAII lock, locks in ctor, unlocks in dtor
//   Use when: you hold the lock for the ENTIRE scope, no early unlock needed
//   (e.g., singleton creation, simple read/write to shared state)
//
// unique_lock<mutex>  -> flexible RAII lock, supports:
//   - early unlock/relock via .unlock()/.lock()
//   - deferred locking (defer_lock)
//   - try_lock (non-blocking attempt)
//   - condition_variable (requires unique_lock, not lock_guard)
//   Use when: you need to unlock mid-scope (e.g., release lock
//   before doing expensive I/O like payment processing)
// ============================================================

class ParkingManager{
    private:
    static ParkingManager* instance;
    static mutex singletonMtx;       // only for singleton creation

    // operationMtx: guards shared mutable state (floors, activeTickets, ticketCounter)
    // WHY mutex here? Multiple threads can call parkVehicle/unparkVehicle concurrently.
    // Without this, two threads could find the same spot available and double-park.
    mutable mutex operationMtx;

    vector<ParkingFloor> floors;
    // shared_ptr<Ticket>: ticket is owned by the map AND returned to caller.
    // Both may need it alive -> shared ownership -> shared_ptr
    map<string, shared_ptr<Ticket>> activeTickets;

    // raw ptr: ParkingManager does NOT own the strategy.
    // Caller owns it and can swap it anytime. Non-owning observer.
    RateStrategy* rateStrategy;
    int ticketCounter;

    ParkingManager() : rateStrategy(nullptr), ticketCounter(0) {}

    public:
    ParkingManager(const ParkingManager&) = delete;
    ParkingManager& operator=(const ParkingManager&) = delete;

    static ParkingManager* getInstance(){
        // Double-checked locking for singleton
        // lock_guard is enough here: we hold the lock for the entire if-block
        if(instance == nullptr){
            lock_guard<mutex> lock(singletonMtx);
            if(instance == nullptr){
                instance = new ParkingManager();
            }
        }
        return instance;
    }

    void setRateStrategy(RateStrategy* strategy){
        lock_guard<mutex> lock(operationMtx);
        rateStrategy = strategy;
    }

    void addFloor(ParkingFloor floor){
        lock_guard<mutex> lock(operationMtx);
        floors.push_back(floor);
    }

    shared_ptr<Ticket> parkVehicle(Vehicle* vehicle){
        // lock_guard: we need the lock for the entire duration of finding a spot
        // and creating a ticket. If we released early, another thread could
        // grab the same spot between our find and park calls (race condition).
        lock_guard<mutex> lock(operationMtx);

        for(auto& floor : floors){
            ParkingSpot* spot = floor.findAvailableSpot(vehicle->getType());
            if(spot != nullptr){
                spot->park(vehicle);
                string ticketId = "TKT-" + to_string(++ticketCounter);
                auto ticket = make_shared<Ticket>(ticketId, *vehicle, spot, chrono::steady_clock::now());
                activeTickets[ticketId] = ticket;
                cout << "[Thread " << this_thread::get_id() << "] "
                     << "Vehicle " << vehicle->getLicensePlate()
                     << " parked at spot " << spot->getSpotId()
                     << " on floor " << floor.getFloorNumber()
                     << " | Ticket: " << ticketId << endl;
                return ticket;
            }
        }
        cout << "[Thread " << this_thread::get_id() << "] "
             << "No available spot for vehicle " << vehicle->getLicensePlate() << endl;
        return nullptr;
    }

    void unparkVehicle(string ticketId, PaymentStrategy* paymentStrategy){
        // unique_lock here instead of lock_guard because:
        // 1. We need the lock to find and update the ticket (shared state)
        // 2. But we can RELEASE it before payment processing (expensive I/O)
        //    so other threads aren't blocked waiting during payment
        // 3. Then re-acquire to erase from map
        unique_lock<mutex> lock(operationMtx);

        auto it = activeTickets.find(ticketId);
        if(it == activeTickets.end()){
            cout << "[Thread " << this_thread::get_id() << "] "
                 << "Ticket " << ticketId << " not found!" << endl;
            return;
        }

        shared_ptr<Ticket> ticket = it->second;
        ticket->setExitTime(chrono::steady_clock::now());
        long long minutes = ticket->getDurationMinutes();
        double amount = rateStrategy ? rateStrategy->calculateRate(minutes) : 0.0;
        ticket->setAmount(amount);
        ticket->getSpot()->unpark();
        ticket->setStatus(TicketStatus::PAID);

        // Release lock BEFORE payment (I/O-heavy, don't block other threads)
        lock.unlock();

        // Payment happens outside the lock - other threads can park/unpark meanwhile
        paymentStrategy->pay(amount);

        // Re-acquire lock to modify shared state (erase from map)
        lock.lock();
        activeTickets.erase(ticketId);

        cout << "[Thread " << this_thread::get_id() << "] "
             << "Vehicle unparked. Ticket " << ticketId << " closed. Amount: $" << amount << endl;
    }

    int getActiveTicketCount(){
        lock_guard<mutex> lock(operationMtx);
        return activeTickets.size();
    }
};

ParkingManager* ParkingManager::instance = nullptr;
mutex ParkingManager::singletonMtx;

#endif
\end{lstlisting}

% ============================================================
\newpage
\section{main.cpp}
% ============================================================
\begin{lstlisting}[title={\texttt{main.cpp}}]
#include<iostream>
#include<thread>
#include<vector>
#include "managers.hpp"
using namespace std;


int main()
{
    ParkingManager* manager = ParkingManager::getInstance();

    // Rate strategies
    HourlyRateStrategy hourlyRate(10.0);
    FlatRateStrategy flatRate(50.0);

    // Set initial rate strategy
    manager->setRateStrategy(&hourlyRate);

    // Create parking spots for floor 1
    vector<ParkingSpot> floor1Spots = {
        ParkingSpot("F1-S1", SpotType::SMALL),
        ParkingSpot("F1-S2", SpotType::MEDIUM),
        ParkingSpot("F1-S3", SpotType::MEDIUM),
        ParkingSpot("F1-S4", SpotType::LARGE)
    };
    ParkingFloor floor1(1, floor1Spots);
    manager->addFloor(floor1);

    // Create parking spots for floor 2
    vector<ParkingSpot> floor2Spots = {
        ParkingSpot("F2-S1", SpotType::SMALL),
        ParkingSpot("F2-S2", SpotType::SMALL),
        ParkingSpot("F2-S3", SpotType::MEDIUM),
        ParkingSpot("F2-S4", SpotType::MEDIUM),
        ParkingSpot("F2-S5", SpotType::LARGE)
    };
    ParkingFloor floor2(2, floor2Spots);
    manager->addFloor(floor2);

    // ============================================================
    // DEMO 1: Concurrent parking from multiple threads
    // Without the lock_guard in parkVehicle, two threads could
    // find the same spot free and both try to park there!
    // ============================================================
    cout << "=== DEMO 1: Concurrent Parking (lock_guard) ===" << endl;
    cout << "Multiple threads trying to park vehicles simultaneously.\n"
         << "lock_guard in parkVehicle ensures only one thread\n"
         << "can find+claim a spot at a time (no double-parking).\n" << endl;

    // Vehicles owned by main thread, passed as raw (non-owning) ptrs to parkVehicle
    Vehicle bike1("BIKE-001", VehicleType::BIKE);
    Vehicle bike2("BIKE-002", VehicleType::BIKE);
    Vehicle car1("CAR-001", VehicleType::CAR);
    Vehicle car2("CAR-002", VehicleType::CAR);
    Vehicle car3("CAR-003", VehicleType::CAR);
    Vehicle truck1("TRUCK-001", VehicleType::TRUCK);
    Vehicle truck2("TRUCK-002", VehicleType::TRUCK);

    // shared_ptr<Ticket>: returned from parkVehicle.
    // Why shared_ptr? The ticket is stored in the manager's map AND
    // we hold a reference here. Both owners -> shared ownership.
    vector<shared_ptr<Ticket>> tickets(7);

    // Spawn threads that concurrently try to park
    vector<thread> parkThreads;
    auto parkFn = [&](int idx, Vehicle* v){
        tickets[idx] = manager->parkVehicle(v);
    };

    parkThreads.emplace_back(parkFn, 0, &bike1);
    parkThreads.emplace_back(parkFn, 1, &bike2);
    parkThreads.emplace_back(parkFn, 2, &car1);
    parkThreads.emplace_back(parkFn, 3, &car2);
    parkThreads.emplace_back(parkFn, 4, &car3);
    parkThreads.emplace_back(parkFn, 5, &truck1);
    parkThreads.emplace_back(parkFn, 6, &truck2);

    for(auto& t : parkThreads) t.join();

    cout << "\nActive tickets after concurrent parking: "
         << manager->getActiveTicketCount() << endl;

    // ============================================================
    // DEMO 2: Concurrent unparking with unique_lock
    // unique_lock lets us RELEASE the lock before payment I/O
    // and RE-ACQUIRE it to erase from the map.
    // This means other threads aren't blocked during payment.
    // ============================================================
    cout << "\n=== DEMO 2: Concurrent Unparking (unique_lock) ===" << endl;
    cout << "unique_lock in unparkVehicle releases the lock before\n"
         << "payment processing, so other threads can work meanwhile.\n"
         << "Then re-acquires to safely erase from the map.\n" << endl;

    // Using HourlyRate for first batch
    CashPayment cashPayment;
    CardPayment cardPayment;

    vector<thread> unparkThreads;

    // Unpark bike1 & car1 with HourlyRate + Cash (threads run concurrently)
    for(int i : {0, 2}){
        if(tickets[i]){
            string tid = tickets[i]->getTicketId();
            unparkThreads.emplace_back([&manager, tid, &cashPayment](){
                manager->unparkVehicle(tid, &cashPayment);
            });
        }
    }

    // Switch to FlatRate mid-flight for the remaining vehicles
    manager->setRateStrategy(&flatRate);

    // Unpark car2 & truck1 with FlatRate + Card (threads run concurrently)
    for(int i : {3, 5}){
        if(tickets[i]){
            string tid = tickets[i]->getTicketId();
            unparkThreads.emplace_back([&manager, tid, &cardPayment](){
                manager->unparkVehicle(tid, &cardPayment);
            });
        }
    }

    for(auto& t : unparkThreads) t.join();

    cout << "\nActive tickets after concurrent unparking: "
         << manager->getActiveTicketCount() << endl;

    // ============================================================
    // SUMMARY OF LOCK CHOICES IN THIS CODEBASE:
    // ----------------------------------------------------------
    // getInstance()     -> lock_guard  (simple scope, create-and-done)
    // setRateStrategy() -> lock_guard  (simple one-liner write)
    // addFloor()        -> lock_guard  (simple one-liner write)
    // getActiveCount()  -> lock_guard  (simple one-liner read)
    // parkVehicle()     -> lock_guard  (must hold lock entire time:
    //                      find spot + park + create ticket is atomic)
    // unparkVehicle()   -> unique_lock (need to unlock before payment
    //                      I/O, then relock to erase from map)
    // ============================================================

    return 0;
}
\end{lstlisting}

\end{document}
